% =====================================================================
%  Section: Method — SceneGraph3D, ViewCtrl, SpatialFlow (v2 additions)
%  These three modules extend the original four-module Canvas3D into
%  Canvas3D-v2, adding viewpoint control and spatial-relation-driven
%  video synthesis on top of the existing perceptual-geometric-
%  compositional-temporal factorization.
% =====================================================================

\subsection{SceneGraph3D: Explicit 3D Scene Graph Construction}
\label{sec:method:scenegraph3d}

\paragraph{Design rationale.}
The original four modules produce a composited 3D scene
$\mathcal{S}^*$ but leave its \emph{relational structure} implicit: the
only relation between two objects is their depth ordering, encoded as a
sort key. This suffices for fixed-camera animation but is insufficient
for two new capabilities we target---\emph{viewpoint-controllable
exploration} and \emph{spatial-relation-anchored video synthesis}---both
of which require reasoning about which object is where, in front of
what, and supporting or occluding which other object. SceneGraph3D
lifts the implicit depth-sorted scene into an \emph{explicit, editable,
serializable 3D scene graph} $\mathcal{G}_\text{scene} = (\mathcal{V},
\mathcal{E})$, whose nodes carry per-object geometric and semantic
attributes and whose edges carry typed, quantified spatial relations.
The graph is the single handshake representation consumed by both
ViewCtrl (\S\ref{sec:method:viewctrl}) and SpatialFlow
(\S\ref{sec:method:spatialflow}), and is the locus of user
\emph{spatial-relation editing}---the user modifies the graph, not the
pixels, and downstream modules re-render and re-generate accordingly.

\paragraph{Nodes and attributes.}
Each foreground object $i$ from PaintParse and a single background
node $b$ become graph nodes. Each node carries:
(i) semantic attributes---category label, salience
$s_i = 0.5\tilde{a}_i + 0.3 c_i + 0.2(1 - \tilde{r}_i)$, combining
normalized mask area $\tilde{a}_i$, detection confidence $c_i$, and
normalized center offset $\tilde{r}_i$;
(ii) geometric attributes---metric 3D center
$\boldsymbol{\mu}_i$, axis-aligned 3D bounding box, volume, height,
and depth range $[d_i^\text{near}, d_i^\text{far}]$ from the
MonoSplat metric rescaling; and
(iii) rendering handles---the 3DGS model path and original 2D mask.
The scene extent
$\boldsymbol{\xi} = (\xi_x, \xi_y, \xi_z)$ is the outer bounding box
of all foreground nodes and defines the spatial scale against which
relation thresholds are normalized, making the graph
scale-invariant.

\paragraph{Typed spatial relations.}
For every ordered pair of foreground nodes $(i, j)$, $i \neq j$,
SceneGraph3D infers up to four typed relations:
\begin{itemize}
    \item \emph{2D projection} (\textsc{left-of},
    \textsc{right-of}, \textsc{above}, \textsc{below}): comparing the
    $x,y$ components of $\boldsymbol{\mu}_i$ and $\boldsymbol{\mu}_j$
    against $\varepsilon_{xy} = 0.1\,\xi_x$, with confidence
    proportional to the normalized offset.
    \item \emph{Depth ordering} (\textsc{in-front-of},
    \textsc{behind}): if $d_i^\text{far} < d_j^\text{near}$ then $i$
    is \textsc{in-front-of} $j$ with confidence
    $(d_j^\text{near} - d_i^\text{far})/\xi_z$; otherwise the
    centroids decide with confidence $0.5$.
    \item \emph{Topological} (\textsc{inside}, \textsc{adjacent},
    \textsc{occludes}): computed from the 2D masks. \textsc{occludes}
    fires when mask overlap exceeds $0.05$ and $i$ is in front of $j$;
    \textsc{inside} when IoU $> 0.9$ and coverage $> 0.8$;
    \textsc{adjacent} when masks are disjoint but within a 5-pixel
    dilation.
    \item \emph{Physical support} (\textsc{supports},
    \textsc{leans-on}): a heuristic that fires when $i$'s bbox bottom
    is within $0.05\,\xi_y$ of $j$'s bbox top and $i$ is above $j$,
    capturing the painterly convention that one object rests on
    another.
\end{itemize}
Each relation edge stores a confidence in $[0,1]$ and, where
applicable, a quantified attribute---3D Euclidean distance, mask
overlap ratio, or depth gap---so that downstream reasoning can be
soft rather than categorical. Every foreground node is additionally
connected to the background node by an \textsc{in-front-of} edge with
confidence $0.95$. The graph is serialized to JSON, making it a
checkpointable, editable artifact.

\paragraph{Scene graph editing.}
The graph exposes a small editing API
(\textsc{move}, \textsc{rotate}, \textsc{scale},
\textsc{add/remove/change-relation}) that SpatialFlow consumes. This is the
key affordance that makes \emph{spatial-relation controllable} video
possible: the user expresses ``move the person to the left of the
building and behind the tree'' as a graph edit, not a text prompt, and
the edit propagates deterministically through ViewCtrl and SpatialFlow
rather than being interpreted stochastically by a language model.

\begin{figure}[t]
    \centering
    \resizebox{\linewidth}{!}{%
    \begin{tikzpicture}[
        font=\scriptsize,
        nodebox/.style={draw=black!55, rounded corners=3pt, thick,
            fill=purple!10, minimum width=2.4cm, minimum height=0.95cm,
            align=center},
        bgbox/.style={nodebox, fill=gray!12},
        edge/.style={->, >=Stealth, thick, draw=black!55},
        reledge/.style={->, >=Stealth, thick, draw=purple!70!black,
            dashed}
    ]
        \node[nodebox] (p) {Person\\{\tiny salience 0.82}};
        \node[nodebox, right=1.4cm of p] (b) {Building\\{\tiny salience 0.65}};
        \node[bgbox, below=1.0cm of $(p)!0.5!(b)$] (bg) {Background};
        \node[nodebox, left=1.4cm of p] (t) {Tree\\{\tiny salience 0.55}};

        \draw[reledge] (p) -- node[above, font=\tiny]
            {\textsc{in-front-of}} (b);
        \draw[reledge] (t) -- node[above, font=\tiny]
            {\textsc{left-of}} (p);
        \draw[reledge] (b) -- node[right, font=\tiny]
            {\textsc{supports}} (bg);
        \draw[reledge] (p) -- node[left, font=\tiny]
            {\textsc{occludes}} (bg);
        \draw[reledge] (t) to[bend right=20] node[below, font=\tiny]
            {\textsc{behind}} (b);
    \end{tikzpicture}}
    \caption{\textbf{SceneGraph3D.} Foreground objects and the
    background become graph nodes carrying geometric and semantic
    attributes; typed, quantified spatial relations---2D projection,
    depth ordering, topological, and physical-support---become edges.
    The graph is editable and serializable, and is the handshake
    representation consumed by ViewCtrl and SpatialFlow.}
    \label{fig:scenegraph3d}
\end{figure}

% =====================================================================
\subsection{ViewCtrl: Viewpoint-Controllable Novel-View Rendering}
\label{sec:method:viewctrl}

\paragraph{Design rationale.}
ArtiFlow in the original Canvas3D renders keyframes along a fixed
cubic-B-spline arc with a small amplitude, which suffices for
parallax-rich ``3D-ish'' animation but does not let the user
\emph{explore} the painting from a chosen viewpoint. ViewCtrl promotes
the camera from a fixed trajectory to a \emph{user-specified
specification}, decoupling viewpoint control from video synthesis: the
user declares where to look from, ViewCtrl renders the multi-view
keyframe stream, and SpatialFlow/ArtiFlow animate between them. Two
design choices follow. First, viewpoint is specified in
\emph{world coordinates against the scene graph}, not in pixel offsets,
so the same trajectory spec portably applies to any painting whose
SceneGraph3D has been built. Second, each rendered frame carries
\emph{per-node visibility masks} recording which foreground objects are
in the frustum and unoccluded, because SpatialFlow's relation embedding
is pooled over the visible node set and must know which nodes
contribute.

\paragraph{Trajectory specification.}
A viewpoint trajectory is a list of camera poses
$\{\mathbf{T}_t\}_{t=1}^{T}$, each pose being a position, a look-at
point, an up vector, and a field of view. Trajectories are generated
in one of five modes:
\emph{manual} (user-supplied poses);
\emph{orbit} (around a focus node at radius $1.5 h_\text{focus}$,
sweeping a user-specified angle);
\emph{dolly} (along the scene main axis, pushing in or pulling out);
\emph{crane} (vertical boom, looking down or up at a fixed target); and
\emph{auto-orbit} (orbit around the most salient foreground node by
default). For orbit, dolly, and crane, the focus node is selected by
$\arg\max_i s_i$ when the user does not specify one, tying viewpoint
control to the SceneGraph3D salience ranking.

\paragraph{Rendering and visibility.}
For each pose $\mathbf{T}_t$, the composited scene $\mathcal{S}^*$
from SketchComp is rendered by differentiable rasterization to produce
an RGB frame $I_t$, a depth map $D_t$, and an edge map
$E_t = \text{Edge}(I_t)$. The per-node visibility mask
$V_t^i \in \{0,1\}$ for foreground node $i$ is computed by projecting
$\boldsymbol{\mu}_i$ into the camera, checking frustum containment,
and comparing the projected depth against $D_t$ at the projected
location: a node is visible iff it lies inside the frustum and the
scene depth at its projection is within a tolerance of the node's
expected depth, indicating that no closer surface occludes it. The
visible node set $\mathcal{V}_t = \{i : V_t^i = 1\}$ is passed to
SpatialFlow as the pooling set for the relation embedding.

\begin{figure}[t]
    \centering
    \resizebox{\linewidth}{!}{%
    \begin{tikzpicture}[
        font=\scriptsize,
        bx/.style={draw=black!55, rounded corners=3pt, thick,
            minimum width=2.0cm, minimum height=0.8cm, align=center},
        block/.style={draw=black!55, rounded corners=3pt, thick,
            fill=cyan!14, minimum width=2.4cm, minimum height=1.0cm,
            align=center},
        ar/.style={->, >=Stealth, thick, draw=black!55}
    ]
        \node[bx, fill=yellow!10] (sg) {SceneGraph $\mathcal{G}_\text{scene}$};
        \node[bx, fill=gray!8, below=0.25cm of sg] (s) {Scene $\mathcal{S}^*$};
        \node[block, right=0.7cm of sg, minimum width=2.6cm] (traj)
            {\textbf{Trajectory spec}\\
             \footnotesize manual/orbit/dolly/\\crane/auto-orbit};
        \node[block, right=0.5cm of traj, minimum width=2.4cm] (rend)
            {\textbf{Multi-view render}\\
             \footnotesize $I_t, D_t, E_t$};
        \node[block, below=0.45cm of rend, minimum width=2.4cm,
              fill=purple!12] (vis)
            {\textbf{Per-node visibility}\\
             \footnotesize $\mathcal{V}_t$};
        \node[bx, fill=gray!8, right=0.5cm of rend, minimum width=2.2cm]
            (out) {Keyframes\\$\{I_t, D_t, E_t, \mathcal{V}_t\}$};

        \draw[ar] (sg) -- (traj);
        \draw[ar] (s) -| (rend);
        \draw[ar] (traj) -- (rend);
        \draw[ar] (rend) -- (out);
        \draw[ar] (rend) -- (vis);
        \draw[ar] (vis) -- (out);
    \end{tikzpicture}}
    \caption{\textbf{ViewCtrl.} A user-specified trajectory
    (manual, orbit, dolly, crane, or auto-orbit) drives multi-view
    rendering of the composited scene, producing RGB, depth, and edge
    keyframes plus per-node visibility sets $\mathcal{V}_t$ that
    determine which scene-graph nodes contribute to SpatialFlow's
    per-frame relation embedding.}
    \label{fig:viewctrl}
\end{figure}

% =====================================================================
\subsection{SpatialFlow: Spatial-Relation-Driven Video Synthesis}
\label{sec:method:spatialflow}

\paragraph{Design rationale.}
ArtiFlow's dual control branches (depth + edge) anchor the video
diffusion backbone to per-frame geometry and boundary, but they carry
no information about \emph{which objects are where relative to each
other}. As a result, the backbone can synthesize frames that are
individually geometrically plausible but collectively violate the
edited scene graph---for example, placing a person in front of a
building in frame $t$ but behind it in frame $t+1$ when the user
edited the graph to specify the opposite. SpatialFlow closes this gap
by injecting a \emph{third control signal}: a GNN-encoded scene-graph
embedding, pooled per-frame over the visible node set, that supplies
the relational context the depth and edge signals cannot. The
embedding is added as a cross-attention context to the video diffusion
transformer blocks, sharing the key/value space with the text prompt,
so that text and graph jointly steer generation. A
relation-consistency verification loop then re-prompts when generated
frames violate the edited graph, closing the loop between editing and
synthesis.

\paragraph{Graph encoding.}
Node features are constructed by concatenating a category one-hot
vector, log-scaled geometric attributes (center, volume, height, depth
range), and the salience score, then linearly projecting to a hidden
dimension $d_h$. Two layers of mean-aggregation message passing
propagate information along the typed relation edges:
\begin{equation}
    \mathbf{h}_i^{(\ell+1)} \;=\;
        \tanh\!\Bigl(\mathbf{h}_i^{(\ell)} \;+\;
        \frac{1}{|\mathcal{N}(i)|}\sum_{j \in \mathcal{N}(i)}
        \mathbf{h}_j^{(\ell)}\Bigr),
    \label{eq:gnn_mp}
\end{equation}
where $\mathcal{N}(i)$ is the set of nodes with an edge to $i$. The
per-frame context vector is the mean of the visible nodes' embeddings:
\begin{equation}
    \mathbf{c}_t^\text{rel} \;=\;
        \frac{1}{|\mathcal{V}_t|}\sum_{i \in \mathcal{V}_t}
        \mathbf{h}_i^{(L)},
    \label{eq:rel_ctx}
\end{equation}
with a fallback to the mean over all foreground nodes when
$\mathcal{V}_t = \emptyset$. This visibility-aware pooling is the key
design choice: as the viewpoint changes, the set of contributing
objects changes, and the relation embedding adapts, preventing the
backbone from receiving stale relational context about objects that
are no longer visible.

\paragraph{Three-way control injection.}
The combined per-frame conditioning signal extends ArtiFlow's dual
control to a triple:
\begin{equation}
    \mathbf{c}_t \;=\;
        s_d\, \mathbf{c}_t^\text{depth}
        \;+\; s_e\, \mathbf{c}_t^\text{edge}
        \;+\; s_r\, \mathbf{c}_t^\text{rel},
    \label{eq:triple_control}
\end{equation}
where $s_d, s_e, s_r$ are the depth, edge, and relation strengths.
The depth and edge branches inject as residual control as in ArtiFlow;
the relation embedding $\mathbf{c}_t^\text{rel}$ is injected as a
cross-attention context in the transformer blocks, sharing the
key/value space with the text-prompt embedding so that text and graph
can jointly steer generation. SpatialFlow supports three generation
modes: \emph{I2V} (first frame = ViewCtrl render + depth + edge +
relation $\to$ video diffusion), \emph{T2V} (text + relation $\to$
first-frame image diffusion $\to$ I2V), and \emph{hybrid} (user
first frame + text + relation $\to$ I2V).

\paragraph{Relation-consistency verification loop.}
After generation, each frame is segmented and the detected object
layout is checked against the edited scene graph: (i) the visibility
of each foreground node is compared against ViewCtrl's
$\mathcal{V}_t$ (visibility IoU); (ii) the edited spatial relations
(e.g.\ ``person left-of building'') are checked against the detected
layout (relation accuracy). If the average visibility IoU falls below
$\tau_\text{vis}$ or any edited relation is violated in more than a
fraction $\tau_\text{rel}$ of frames, SpatialFlow triggers a
\emph{re-prompt}: it increases $s_r$, optionally appends an explicit
textual description of the violated relation to the prompt, and
regenerates the affected segment. This closes the loop between scene
graph editing and video synthesis, and is the direct generalization of
ArtiFlow's re-composition loop to the relational setting.

\begin{figure}[t]
    \centering
    \resizebox{\linewidth}{!}{%
    \begin{tikzpicture}[
        font=\scriptsize,
        bx/.style={draw=black!55, rounded corners=3pt, thick,
            minimum width=2.0cm, minimum height=0.8cm, align=center},
        block/.style={draw=black!55, rounded corners=3pt, thick,
            fill=red!12, minimum width=2.4cm, minimum height=1.0cm,
            align=center},
        gnblock/.style={draw=black!55, rounded corners=3pt, thick,
            fill=purple!14, minimum width=2.4cm, minimum height=1.0cm,
            align=center},
        ar/.style={->, >=Stealth, thick, draw=black!55},
        feedback/.style={->, >=Stealth, thick, draw=red!65, dashed}
    ]
        \node[bx, fill=yellow!10] (sg) {$\mathcal{G}_\text{scene}$\\(edited)};
        \node[gnblock, right=0.6cm of sg] (gnn)
            {\textbf{GNN encoder}\\
             \footnotesize Eq.~\ref{eq:gnn_mp}};
        \node[bx, fill=purple!10, below=0.3cm of gnn, minimum width=2.2cm]
            (ctx) {$\mathbf{c}_t^\text{rel}$\\(Eq.~\ref{eq:rel_ctx})};
        \node[bx, fill=gray!8, right=0.6cm of gnn, minimum width=2.0cm]
            (kf) {ViewCtrl\\keyframes};
        \node[block, below=0.6cm of kf, minimum width=3.0cm] (cn)
            {\textbf{Triple control}\\
             \footnotesize depth + edge + relation\\
             \footnotesize (Eq.~\ref{eq:triple_control})};
        \node[block, right=0.5cm of cn, minimum width=2.6cm,
              minimum height=1.4cm] (vd)
            {\textbf{Video diffusion}\\
             \footnotesize cross-attn relation\\
             \footnotesize + first-frame anchor};
        \node[bx, fill=red!18, below=0.4cm of vd, minimum width=2.4cm]
            (out) {Video $V$};
        \node[gnblock, below=0.55cm of out, minimum width=2.8cm]
            (verify) {\textbf{Relation verify}\\
             \footnotesize visibility IoU +\\relation accuracy};

        \draw[ar] (sg) -- (gnn);
        \draw[ar] (gnn) -- (ctx);
        \draw[ar] (kf) -- (cn);
        \draw[ar] (ctx) -| (cn);
        \draw[ar] (cn) -- (vd);
        \draw[ar] (kf) |- (vd);
        \draw[ar] (vd) -- (out);
        \draw[ar] (out) -- (verify);
        \draw[feedback] (verify.west) -- ++(-1.3,0)
            node[midway, above, font=\tiny, text=red!75]
            {re-prompt} |- (vd.west);
    \end{tikzpicture}}
    \caption{\textbf{SpatialFlow.} The edited scene graph is encoded
    by a GNN and pooled per-frame over the visible node set into a
    relation context $\mathbf{c}_t^\text{rel}$, which is injected as a
    third control signal alongside depth and edge into the video
    diffusion backbone. A relation-consistency verification loop
    re-prompts when generated frames violate the edited graph.}
    \label{fig:spatialflow}
\end{figure}
